\chapter{Mathematics}
\label{chap:mathematics}

\section{Equations}

The Gaussian integral is a classic:
\begin{equation}
  \int_{-\infty}^{\infty} e^{-x^2}\,\mathrm{d}x = \sqrt{\pi}.
  \label{eq:gauss}
\end{equation}
Aligned derivations keep every step readable. Here is the least squares
solution of $\min_{\beta} \norm{y - X\beta}^2$:
\begin{align}
  \hat\beta &= \argmin_{\beta \in \R^p} \norm{y - X\beta}_2^2 \label{eq:ls} \\
            &= \left(X^\top X\right)^{-1} X^\top y. \nonumber
\end{align}

Piecewise functions and matrices are just as easy:
\[
  f(x) = \begin{cases}
    x^2        & \text{if } x \ge 0,\\
    -\sin x    & \text{otherwise,}
  \end{cases}
  \qquad
  R_\theta = \begin{pmatrix}
    \cos\theta & -\sin\theta \\
    \sin\theta & \cos\theta
  \end{pmatrix}.
\]

\section{Theorems}

\begin{definition}[Metric space]
  A metric space is a pair $(M, d)$ where $d \colon M \times M \to \R$ is
  non-negative, symmetric, zero exactly on the diagonal, and satisfies the
  triangle inequality.
\end{definition}

\begin{theorem}[Banach fixed point]
  \label{thm:banach}
  Let $(M,d)$ be a non-empty complete metric space and $T\colon M\to M$ a
  contraction. Then $T$ has a unique fixed point $x^\ast$, and for every
  $x_0 \in M$ the sequence $x_{n+1} = T(x_n)$ converges to $x^\ast$.
\end{theorem}

\begin{proof}
  The iterates form a Cauchy sequence, because
  $d(x_{n+1},x_n) \le q^n d(x_1,x_0)$ with $q<1$. Completeness gives a limit,
  continuity of $T$ makes it a fixed point, and contractivity makes it unique.
\end{proof}

\begin{remark}
  \Cref{thm:banach} and \cref{eq:gauss,eq:ls} are clickable links in the PDF.
\end{remark}

\section{Units}

With \texttt{siunitx}, quantities are typeset consistently: the speed of light
is \qty{299792458}{\metre\per\second}, and a laptop GPU draws about
\qty{60}{\watt} at \qty{1.5}{\giga\hertz}.
